\chapter{Deudas y trabajo futuro}
\label{sec:deudas}

Este capítulo documenta lo que \emph{no} está resuelto. Hay cuatro
categorías: aristas indebidas en el grafo de dependencias, tests
flaky, brechas de modelo y follow-ups del sistema agentic.

\section{Dos aristas \filepath{core} $\to$ \filepath{platform}
(con \emph{guardrail} automatizado)}

Como registra el capítulo~\ref{sec:grafo-deps}, existen dos
dependencias que van hacia arriba:

\begin{itemize}
    \item \filepath{core/graph-engine} $\to$
    \filepath{platform/sources}: el GNN bridge usa
    \code{AadEventBatch} definido en \filepath{platform/sources}
    para materializar el corpus de entrenamiento. Es un préstamo de
    tipo; se resuelve moviendo \code{AadEventBatch} a un crate más
    bajo (candidato: un crate nuevo \filepath{core/ingest-types})
    o invirtiendo la dirección vía un trait.
    \item \filepath{core/graph-engine} $\to$ \filepath{platform/dsl}:
    el compilador de hipótesis y el runtime del grafo comparten
    tipos; separarlos completamente exige introducir un crate
    intermedio de tipos compartidos.
\end{itemize}

No son ciclos (el grafo sigue siendo DAG) pero violan la regla de
capas. \textbf{A partir del ciclo 2026-04 ambas aristas están
explícitamente \emph{whitelisted} en
\filepath{docs/arquitectura/layering-exceptions.yaml} y verificadas
en cada PR por \filepath{scripts/check-layering.sh}}. Esto convierte
la deuda en regresión visible: cualquier nueva arista
\filepath{core} $\to$ \filepath{platform} no-whitelisted rompe CI.
El indicador de éxito de largo plazo sigue siendo que
\code{cargo metadata} reporte cero aristas hacia arriba; el
indicador de corto plazo ---cerrado--- es que ninguna arista nueva
pueda colarse silenciosamente.

\section{Tests flaky conocidos}

\begin{itemize}
    \item \code{sentinel\_connector::tests::backoff\_retries} —
    depende de timing con \code{tokio::time::advance}, a veces
    falla en CI bajo carga. Mitigación actual: reintentar el job.
    Solución: rediseñar el test para usar relojes virtuales sin
    depender de wall-clock.
    \item Smoke tests de ingesta que abren archivos grandes en
    paralelo saturan I/O en runners GitHub y producen timeouts
    esporádicos. Mitigación: serializar la fase de ingesta en CI.
\end{itemize}

Ninguno se ha visto fallar localmente de forma reproducible. No
bloquean releases pero contaminan el ruido de CI.

\section{GNN v2 sin modelo entrenado}

ADR-003 previene que un binario acepte el modelo skeleton, pero no
resuelve la brecha subyacente: el entrenamiento de GNN v2 está
hecho sobre Sentinel/AAD y no se ha validado contra datasets
Fortigate o IIS. Dos consecuencias:

\begin{itemize}
    \item La componente \code{gnn\_threat} del scoring compuesto no
    aporta señal útil cuando se alimenta con logs fuera del dominio
    entrenado. En esos casos el gate rechaza la carga y el scoring
    se reduce a cuatro componentes. Es el comportamiento correcto,
    pero es menos discriminante.
    \item No hay \emph{corpus export} operativo para
    FortiAnalyzer; el trabajo necesario incluye estructurar el
    pipeline de exportación + entrenamiento para esa familia de
    logs.
\end{itemize}

\section{Follow-ups agentic slow lane}

El sistema agentic (integrado en Fase 2 vía ADR-006) tiene tres
piezas pendientes:

\begin{itemize}
    \item \textbf{Backend local-llm real}: actualmente el backend
    \code{candle} está cableado pero el modelo de bajo recurso no
    converge con calidad suficiente para el \emph{review loop}.
    Falta decidir entre (a) afinar el modelo actual, (b) integrar
    un backend remoto con API key rotada, o (c) ambos con fallback.
    \item \textbf{Drift detection feedback loop}: la operación
    \code{agentic::drift} detecta schema drift pero no propaga la
    corrección al \code{mapping\_library}. Falta cablear esa vuelta
    para que un drift detectado genere una propuesta de re-mapping
    revisable.
    \item \textbf{DLQ triage UI}: la DLQ es funcional por API pero
    no tiene vista en el Tauri. Consecuencia: un operador ve
    ``hay N eventos fallidos'' pero no puede triagearlos sin
    \code{curl}.
\end{itemize}

\section{Deudas cerradas en el ciclo 2026-04}

Se dejan registradas aquí, con puntero a la evidencia, para que la
próxima revisión del capítulo no tenga que re-descubrirlas:

\begin{itemize}
    \item \textbf{FFI sin anotaciones \code{SAFETY}}: todos los
    bloques \code{unsafe} de \filepath{core/matcher-ffi/src/lib.rs}
    tienen ahora un comentario \code{// SAFETY: ...} explícito y el
    lint \code{missing-safety-doc} está habilitado. Test de ABI
    mismatch + 6 tests negativos (null, grafo vacío, patrón
    inválido, ventana temporal fuera de rango, double-drop bloqueado
    por typestate) en \filepath{core/matcher-ffi/tests/}.

    \item \textbf{DSL v1 sin semántica formal}: publicada en
    \filepath{docs/spec/dsl-v1-semantics.md}. Cubre definición de
    \emph{match} (paths vs walks), semántica de \code{WITHIN},
    \code{\{k=N\}}, \code{HAVING count(distinct)}, y los predicados
    \code{\textasciitilde}, \code{in}, \code{not in}. El roundtrip
    proptest cubre ahora bindings, HAVING y WITHIN que antes se
    excluían.

    \item \textbf{Matcher 3-way equivalence}: proptest en
    \filepath{core/matcher-ffi/tests/equivalence\_proptest.rs}
    compara iterativo vs smart vs SIMD a 256 casos. Con scorer
    finalizado, los tres devuelven conjuntos idénticos.

    \item \textbf{Scoring sin \code{ScoreBreakdown} documentado}:
    publicado en \filepath{docs/spec/scoring-v1.md}. Perfil por
    defecto con \code{w\_gnn = 0}, breakdown por path expuesto en
    \code{analytics::HuntResult::anomaly\_breakdown}, tests de
    reproducibilidad bit-for-bit (\code{f64::to\_bits}) y
    direccionalidad.

    \item \textbf{Escenarios demo no reproducibles}:
    \filepath{scripts/demo.sh} + snapshot tests en
    \filepath{platform/api/tests/demo\_scenarios.rs} cubren los
    tres fixtures (Sysmon, Sentinel, CSV) con tolerancia
    $\pm 10\%$ sobre el número de paths.

    \item \textbf{Benchmarks limitados a 2 suites}: agregadas
    \code{ingest\_throughput}, \code{matching\_hops} y
    \code{scoring\_path}. \filepath{scripts/bench-report.sh}
    convierte \filepath{target/criterion/*/estimates.json} en
    \filepath{docs/paper/data/bench-table.md}.

    \item \textbf{Feature flags sin tabla canónica}: publicada en
    \filepath{docs/spec/features.md} con clasificación
    stable/experimental/future sobre 14 flags.
\end{itemize}

\section{Formato del libro}

Este capítulo se actualiza en cada ciclo de trabajo. Si una deuda
se resuelve, se elimina de aquí y se anota en la ADR correspondiente
bajo ``consecuencias''. Si aparece una nueva, se documenta aquí
antes de que alguien la descubra por sorpresa durante un incidente.

La regla editorial: las deudas se enumeran con nombre de módulo,
efecto observable y criterio de resolución. No se usan frases
vagas (``mejorar''); se usan criterios verificables (``cero
aristas hacia arriba en \code{cargo metadata}'').
